% ===========================================================================
% 15. Conclusion
% ===========================================================================
\section{Conclusion}
\label{sec:conclusion}

We presented \textsc{agent-traps-lab}, an open-source testbed that
bridges the gap between the theoretical DeepMind AI Agent Traps
taxonomy and empirical evaluation. By operationalizing all six trap
categories as 22 reproducible adversarial scenarios and evaluating
them against four production-grade LLMs over 9{,}120 controlled
experiments, we provide the first comprehensive, statistically
rigorous measurement of environmental attacks on multi-agent A2A
systems.

Our key findings are:

\begin{enumerate}[nosep,leftmargin=*]
  \item \textbf{All six trap categories are empirically effective}
    against undefended agents, with content injection achieving the
    highest baseline success rates and semantic manipulation being
    the hardest to detect.

  \item \textbf{Model vulnerability varies significantly}: efficient
    models (\model{GPT-4o-mini}) are consistently more vulnerable
    than frontier models, but all models exhibit category-specific
    weaknesses.

  \item \textbf{Defense-in-depth works}: the seven-module mitigation
    suite reduces overall trap success by a statistically significant
    margin with large effect sizes across all models and categories.

  \item \textbf{Compound traps are super-additive}: combining two
    trap categories produces success rates higher than either
    individual category, with content injection + semantic manipulation
    being the most potent combination.

  \item \textbf{Ablation reveals critical modules}: input sanitization
    and cascade breaking provide the greatest marginal defense value,
    while RAG integrity is a single point of failure for cognitive
    state attacks.

  \item \textbf{The A2A protocol needs security extensions}: our
    systemic attack results motivate mandatory message authentication,
    content integrity verification, and trust scoring in future
    protocol versions.
\end{enumerate}

We release the complete testbed---all 22 adversarial scenarios,
7 mitigation modules, experiment harness, and statistical analysis
scripts---as open source (\cref{app:artifacts}) to enable
reproducible adversarial evaluation of multi-agent systems.

\paragraph{Future work.} We plan to extend the testbed with
adaptive adversaries that adjust their strategies based on the
target model's responses, a human-subjects study to validate
the HL-category results, and integration with additional
multi-agent protocols beyond A2A.
